%--------------------------------------------------
% Function: \subinput
%
% Desc:
%   Inputs the contents of another TeX file and 
%   temporarily redefines \section to behave as 
%   \subsection, allowing nested headings in a 
%   larger document.
%
% Arguments:
%   #1 : Path to the TeX file to input. Must be a 
%        valid file relative to the current working 
%        directory.
%
% Example:
%
% \subinput{my_tex_file}
%--------------------------------------------------
\makeatletter
\newcommand{\subinput}[1]{%
  \begingroup
    % Change \section to \subsection
    \let\section\subsection  
    \let\subsection\subsubsection  
    % Input the specified file
    \input{#1}%               
  \endgroup
}
\makeatother

%--------------------------------------------------
% Function: \makeheader
%
% Desc:
%   Sets the master-document flag *and* creates a
%   left-aligned document header with optional
%   description, title, author, date, table of
%   contents, and a page break. The title is also
%   added as a hidden section for referencing.
%
% Arguments:
%   #1 : Optional description text.  Omit to skip.
%   #2 : Required title.
%--------------------------------------------------
\NewDocumentCommand{\makeheader}{ o m }{%
  \masteron % <-- set the master flag
  % Hidden section for the title (TOC + label)
  \section*{#2}%
  \label{first}%
  % Large bold title text
  \begin{flushleft}
    \vspace{0.75em}%
    \textbf{Author:} Jacob S. Zelko\par% Author
    \textbf{Date: }\today\par% Date
    \IfNoValueF{#1}{%
      \small\textbf{Description:}~#1\par
    }%
  \end{flushleft}
  \tableofcontents
  \clearpage
}



% --------------------------------------------------
% Master/child detection
% --------------------------------------------------
\newif\ifmasterdoc       % create a boolean we can set in the parent
\newcommand{\masteron}{\masterdoctrue}
\newcommand{\masteroff}{\masterdocfalse}

% --------------------------------------------------
% Convenience macro to print bibliography only
% when *not* building the master document.
% --------------------------------------------------
\newcommand{\printbib}{%
  \ifmasterdoc\else
    \printbibliography
  \fi
}

%--------------------------------------------------
% Function: \keywords
%
% Desc:
%   Prints a “Keywords” section flush-left, using
%   the same font/size as a \section heading, and
%   adds an entry to the table of contents.
%
% Arguments:
%   #1 : A comma-separated list or free string of
%        keywords to display.
%
% Example:
%   \keywords{data science, category theory, OMOP}
%--------------------------------------------------
\newcommand{\keywords}[1]{%
  % Left-aligned block styled like a section header
  \begin{flushleft}
    {\normalfont\bfseries Keywords:} {\normalfont #1}\par
  \end{flushleft}%
}
